% GENERATED FILE. Edit canonical lessons, then run npm run generate:books.
% canonical-source-hash: fnv1a64:1381dcc4fc4f0b64
% canonical-lessons: JA-C508-un, JA-C508-kouun, JA-C508-sainan, JA-C508-un-yoku, JA-C508-unmei

\chapter{Good Luck and Bad Luck}
\label{ch:ja-good-luck-and-bad-luck}
\hlchaptermodality{508}

\begin{chapteropening}
\textbf{By the end of this chapter:} \emph{I can say what might happen with a little luck.}

Complete the last lesson of chapter 508: I can say what might happen with a little luck.
\end{chapteropening}

\section[\texorpdfstring{\ja{うん} (un)}{うん (un)}]{\ja{うん} (un) — luck; fortune}

\label{lesson:JA-C508-un}



\begin{quote}
Before the new one: say the Japanese for an ideal, then the Japanese for to want (said of someone else).
\end{quote}

\subsection*{You'll want to know: \ja{うん}}
\textbf{\ja{うん}} — \emph{un} — \textquotedblleft{}luck; fortune\textquotedblright{}.

\begin{quote}
\textbf{\ja{もし} \ja{うんが} \ja{よかったら、あには} \ja{しあいに} \ja{かつでしょう。}} — \emph{moshi un ga yokattara, ani wa shiai ni katsu deshō} — \textquotedblleft{}If he is lucky, my elder brother will probably win the match.\textquotedblright{}
\end{quote}

\begin{cousinweb}
\textbf{\ja{うん}} (\emph{un}) is Sino-Japanese (kango): one Chinese character read with its on'yomi, reading borrowed from Chinese. It is \emph{un}, \textquotedblleft{}luck\textquotedblright{} (Middle Chinese \emph{hjunH}) (KANJIDIC2; Middle Chinese in Baxter's transcription).
\end{cousinweb}

\begin{grammarlens}[title={What you've built}]
One more word for saying what happens if.
\end{grammarlens}

\subsection*{Your turn}
\begin{itemize}
\raggedright
  \item \emph{Say it:} \emph{un}
  \item \emph{Say it:} \emph{un}, once more
  \item \emph{Say it:} \emph{moshi un ga yokattara, ani wa shiai ni katsu deshō} — If he is lucky, my elder brother will probably win the match.
  \item \emph{Recall:} say the Japanese for an ideal, then the Japanese for to want (said of someone else), then say \emph{un} again
\end{itemize}

\begin{tcolorbox}[breakable,colback=teal!4,colframe=teal!35!black,title={Before you move on}]
What does \ja{うん} mean? (\textquotedblleft{}Luck; fortune\textquotedblright{}.) Say it once more.
\end{tcolorbox}
\section[\texorpdfstring{\ja{こううん} (kōun)}{こううん (kōun)}]{\ja{こううん} (kōun) — good fortune}

\label{lesson:JA-C508-kouun}



\begin{quote}
Before the new one: say the Japanese for to want (said of someone else), then the Japanese for luck.
\end{quote}

\subsection*{You'll want to know: \ja{こううん}}
\textbf{\ja{こううん}} — \emph{kōun} — \textquotedblleft{}good fortune\textquotedblright{}.

\begin{quote}
\textbf{\ja{もし} \ja{いい} \ja{ともだちに} \ja{あったら、それは} \ja{こううんでしょう。}} — \emph{moshi ii tomodachi ni attara, sore wa kōun deshō} — \textquotedblleft{}If you met a good friend, that would be good fortune.\textquotedblright{}
\end{quote}

\begin{cousinweb}
\textbf{\ja{こううん}} (\emph{kōun}) is Sino-Japanese (kango): two Chinese characters read with their on'yomi, readings borrowed from Chinese. They are \emph{kō}, \textquotedblleft{}happiness\textquotedblright{} (Middle Chinese \emph{heangX}) and \emph{un}, \textquotedblleft{}luck\textquotedblright{} (\emph{hjunH}) (KANJIDIC2; Middle Chinese in Baxter's transcription).
\end{cousinweb}

\begin{grammarlens}[title={What you've built}]
One more word for saying what happens if.
\end{grammarlens}

\subsection*{Your turn}
\begin{itemize}
\raggedright
  \item \emph{Say it:} \emph{kōun}
  \item \emph{Say it:} \emph{kōun}, once more
  \item \emph{Say it:} \emph{moshi ii tomodachi ni attara, sore wa kōun deshō} — If you met a good friend, that would be good fortune.
  \item \emph{Recall:} say the Japanese for to want (said of someone else), then the Japanese for luck, then say \emph{kōun} again
\end{itemize}

\begin{tcolorbox}[breakable,colback=teal!4,colframe=teal!35!black,title={Before you move on}]
What does \ja{こううん} mean? (\textquotedblleft{}Good fortune\textquotedblright{}.) Say it once more.
\end{tcolorbox}
\section[\texorpdfstring{\ja{さいなん} (sainan)}{さいなん (sainan)}]{\ja{さいなん} (sainan) — a misfortune; a calamity}

\label{lesson:JA-C508-sainan}



\begin{quote}
Before the new one: say the Japanese for luck, then the Japanese for good fortune.
\end{quote}

\subsection*{You'll want to know: \ja{さいなん}}
\textbf{\ja{さいなん}} — \emph{sainan} — \textquotedblleft{}a misfortune; a calamity\textquotedblright{}.

\begin{quote}
\textbf{\ja{もし} \ja{りょこうの} \ja{まえに} \ja{けがを} \ja{したら、さいなんですね。}} — \emph{moshi ryokō no mae ni kega o shitara, sainan desu ne} — \textquotedblleft{}If you got hurt right before a trip, that would be bad luck.\textquotedblright{}
\end{quote}

\begin{cousinweb}
\textbf{\ja{さいなん}} (\emph{sainan}) is Sino-Japanese (kango): two Chinese characters read with their on'yomi, readings borrowed from Chinese. They are \emph{sai}, \textquotedblleft{}disaster\textquotedblright{} (Middle Chinese \emph{tsoj}) and \emph{nan}, \textquotedblleft{}trouble\textquotedblright{} (\emph{nanH}) (KANJIDIC2; Middle Chinese in Baxter's transcription).
\end{cousinweb}

\begin{grammarlens}[title={What you've built}]
One more word for saying what happens if.
\end{grammarlens}

\subsection*{Your turn}
\begin{itemize}
\raggedright
  \item \emph{Say it:} \emph{sainan}
  \item \emph{Say it:} \emph{sainan}, once more
  \item \emph{Say it:} \emph{moshi ryokō no mae ni kega o shitara, sainan desu ne} — If you got hurt right before a trip, that would be bad luck.
  \item \emph{Recall:} say the Japanese for luck, then the Japanese for good fortune, then say \emph{sainan} again
\end{itemize}

\begin{tcolorbox}[breakable,colback=teal!4,colframe=teal!35!black,title={Before you move on}]
What does \ja{さいなん} mean? (\textquotedblleft{}A misfortune; a calamity\textquotedblright{}.) Say it once more.
\end{tcolorbox}
\section[\texorpdfstring{\ja{うんよく} (un'yoku)}{うんよく (un'yoku)}]{\ja{うんよく} (un'yoku) — luckily}

\label{lesson:JA-C508-un-yoku}



\begin{quote}
Before the new one: say the Japanese for good fortune, then the Japanese for a misfortune.
\end{quote}

\subsection*{You'll want to know: \ja{うんよく}}
\textbf{\ja{うんよく}} — \emph{un'yoku} — \textquotedblleft{}luckily\textquotedblright{}.

\begin{quote}
\textbf{\ja{うんよく} \ja{でんしゃに} \ja{まにあったら、ちちは} \ja{ごじに} \ja{つくでしょう。}} — \emph{un'yoku densha ni maniattara, chichi wa goji ni tsuku deshō} — \textquotedblleft{}If my father is lucky and catches the train, he will arrive at five.\textquotedblright{}
\end{quote}

\begin{cousinweb}
\textbf{\ja{うんよく}} (\emph{un'yoku}) mixes the two layers, native and Sino-Japanese: Sino-Japanese \emph{un}, \textquotedblleft{}luck\textquotedblright{} (Middle Chinese \emph{hjunH}), borrowed from Chinese, and native \emph{yoku}, \textquotedblleft{}well\textquotedblright{}, the adverb of \emph{yoi}, \textquotedblleft{}good\textquotedblright{} (JMdict, KANJIDIC2; Middle Chinese in Baxter's transcription).
\end{cousinweb}

\begin{grammarlens}[title={What you've built}]
One more word for saying what happens if.
\end{grammarlens}

\subsection*{Your turn}
\begin{itemize}
\raggedright
  \item \emph{Say it:} \emph{un'yoku}
  \item \emph{Say it:} \emph{un'yoku}, once more
  \item \emph{Say it:} \emph{un'yoku densha ni maniattara, chichi wa goji ni tsuku deshō} — If my father is lucky and catches the train, he will arrive at five.
  \item \emph{Recall:} say the Japanese for good fortune, then the Japanese for a misfortune, then say \emph{un'yoku} again
\end{itemize}

\begin{tcolorbox}[breakable,colback=teal!4,colframe=teal!35!black,title={Before you move on}]
What does \ja{うんよく} mean? (\textquotedblleft{}Luckily\textquotedblright{}.) Say it once more.
\end{tcolorbox}
\section[\texorpdfstring{\ja{うんめい} (unmei)}{うんめい (unmei)}]{\ja{うんめい} (unmei) — fate; destiny}

\label{lesson:JA-C508-unmei}



\begin{quote}
Before the new one: say the Japanese for a misfortune, then the Japanese for luckily.
\end{quote}

\subsection*{You'll want to know: \ja{うんめい}}
\textbf{\ja{うんめい}} — \emph{unmei} — \textquotedblleft{}fate; destiny\textquotedblright{}.

\begin{quote}
\textbf{\ja{もし} \ja{あねが} \ja{まちで} \ja{かれに} \ja{あったら、うんめいでしょう。}} — \emph{moshi ane ga machi de kare ni attara, unmei deshō} — \textquotedblleft{}If my elder sister met him in town, it would be fate.\textquotedblright{}
\end{quote}

\begin{cousinweb}
\textbf{\ja{うんめい}} (\emph{unmei}) is Sino-Japanese (kango): two Chinese characters read with their on'yomi, readings borrowed from Chinese. They are \emph{un}, \textquotedblleft{}luck\textquotedblright{} (Middle Chinese \emph{hjunH}) and \emph{mei}, \textquotedblleft{}fate\textquotedblright{} (\emph{mjaengH}) (KANJIDIC2; Middle Chinese in Baxter's transcription).
\end{cousinweb}

\begin{grammarlens}[title={What you've built}]
One more word for saying what happens if.
\end{grammarlens}

\subsection*{Your turn}
\begin{itemize}
\raggedright
  \item \emph{Say it:} \emph{unmei}
  \item \emph{Say it:} \emph{unmei}, once more
  \item \emph{Say it:} \emph{moshi ane ga machi de kare ni attara, unmei deshō} — If my elder sister met him in town, it would be fate.
  \item \emph{Recall:} say the Japanese for a misfortune, then the Japanese for luckily, then say \emph{unmei} again
\end{itemize}

\begin{tcolorbox}[breakable,colback=teal!4,colframe=teal!35!black,title={Before you move on}]
What does \ja{うんめい} mean? (\textquotedblleft{}Fate; destiny\textquotedblright{}.) Say it once more.
\end{tcolorbox}
